% Part: propositional-logic
% Chapter: syntax-and-semantics
% Section: semantic-notions

\documentclass[../../../include/open-logic-section]{subfiles}

\begin{document}

\olfileid{pl}{syn}{sem}

\olsection{चिन्हार्थविषयक संकल्पना}

पुढील चिन्हार्थविषयक संकल्पना परिभाषित करू:

\begin{defn} 
\begin{enumerate}
\item !!^a{formula}~$!A$ हे \emph{पूर्ततायोग्य} असते, जर एखाद्या
  $\pAssign{v}$ साठी $\pSat{v}{!A}$; आणि कोणत्याही $\pAssign{v}$ साठी
  $\pSat{v}{!A}$ नसल्यास ते \emph{अपूर्ततायोग्य} असते;
\item सर्व !!{valuation}s~$\pAssign{v}$ साठी $\pSat{v}{!A}$ असेल, तर
  !!^a{formula}~$!A$ हे \emph{सर्वतः सत्य सूत्र} असते;
\item !!^a{formula}~$!A$ हे पूर्ततायोग्य असेल, पण सर्वतः सत्य सूत्र नसेल,
  तर ते \emph{परिस्थितीवश} असते;
\item $\Gamma$ हा !!{formula}s चा संच असेल, तर $\Gamma \Entails !A$
  (``$\Gamma$ पासून $!A$ तार्किकरीत्या निष्पन्न होते'') तेव्हा आणि केवळ
  तेव्हाच, जेव्हा $\pSat{v}{\Gamma}$ असलेल्या प्रत्येक
  !!{valuation}~$\pAssign{v}$ साठी $\pSat{v}{!A}$.
\item $\Gamma$ हा !!{formula}s चा संच असेल, तर एखाद्या
  !!a{valuation}~$\pAssign{v}$ साठी $\pSat{v}{\Gamma}$ असल्यास $\Gamma$
  हा \emph{पूर्ततायोग्य} असतो; अन्यथा $\Gamma$ हा
  \emph{अपूर्ततायोग्य} असतो.
\end{enumerate} 
\end{defn}

\begin{prob} 
 पुढील चार !!{formula}s पैकी प्रत्येक सूत्र (a)~पूर्ततायोग्य,
 (b)~सर्वतः सत्य सूत्र आणि (c)~परिस्थितीवश आहे की नाही ते ठरवा.
\begin{enumerate}
  \item \( ( \Obj p_0 \lif ( \lnot \Obj p_1 \lif \lnot \Obj p_0 ) ) \).
  \item \( ( ( \Obj p_0 \land \lnot \Obj p_1 ) \lif ( \lnot \Obj p_0 \land \Obj p_2 )) \liff ( ( \Obj p_2 \lif \Obj p_0 ) \lif ( \Obj p_0 \lif \Obj p_1 )) \).
  \item \( ( \Obj p_0 \liff \Obj p_1 ) \lif ( \Obj p_2 \liff \lnot \Obj p_1 ) \).
  \item \( (( \Obj p_0 \liff ( \lnot \Obj p_1 \land \Obj p_2 )) \lor ( \Obj p_2 \lif ( \Obj p_0 \liff \Obj p_1 ))) \).
\end{enumerate}
\end{prob}

\begin{prop}
\ollabel{prop:semanticalfacts} 
\begin{enumerate} 
\item $!A$ हे सर्वतः सत्य सूत्र तेव्हा आणि केवळ तेव्हाच असते, जेव्हा
  $\emptyset \Entails !A$; 
\item $\Gamma \Entails !A$ आणि $\Gamma \Entails !A \lif !B$ असल्यास
  $\Gamma \Entails !B$;
\item $\Gamma$ पूर्ततायोग्य असेल, तर $\Gamma$ चा प्रत्येक सांत उपसंचही
  पूर्ततायोग्य असतो; 
\item \ollabel{def:monotonicity} एकस्वनिकता: $\Gamma \subseteq \Delta$
  आणि $\Gamma \Entails !A$ असल्यास $\Delta \Entails !A$ हेही असते;
\item \ollabel{def:Cut} संक्रमकता: $\Gamma \Entails !A$ आणि
  $\Delta \cup \{ !A\} \Entails !B$ असल्यास $\Gamma \cup \Delta \Entails
  !B$.
\end{enumerate}
\end{prop}

\begin{proof}
सरावासाठी.
\end{proof}

\begin{prob}
\olref[pl][syn][sem]{prop:semanticalfacts} सिद्ध करा.
\end{prob}

\begin{prop}\ollabel{prop:entails-unsat}
  $\Gamma \Entails !A$ तेव्हा आणि केवळ तेव्हाच, जेव्हा
  $\Gamma \cup \{\lnot !A\}$ अपूर्ततायोग्य असतो.
\end{prop}

\begin{proof}
सरावासाठी.
\end{proof}

\begin{prob}
\olref[pl][syn][sem]{prop:entails-unsat} सिद्ध करा.
\end{prob}

\begin{thm}[चिन्हार्थविषयक निगमन प्रमेय]
  \ollabel{thm:sem-deduction} $\Gamma \Entails !A \lif !B$ तेव्हा आणि केवळ
  तेव्हाच, जेव्हा $\Gamma \cup \{!A\} \Entails !B$.
\end{thm}

\begin{proof}
सरावासाठी.
\end{proof}

\begin{prob}
\olref[pl][syn][sem]{thm:sem-deduction} सिद्ध करा.
\end{prob}

\end{document}
